% ============================================================================
% DT-GSK notation table -- CORE SYMBOLS (main text, Section 3.1).
% Canonical source: papers/build_prompt_phases/phase_03/notation_table.md
% MDPI-compatible: article class, booktabs; no empirical values.
% Label: tab:notation (exhibit T-NOTATION, phase_04/exhibit_plan.csv).
%
% Split under ticket N-015: the combined ~49-row key mixed problem variables,
% GSK operators, and every subsystem's internal constants into one dense
% \scriptsize page that readers met before the architecture figure. This file
% now carries only the problem/population and base-GSK symbols needed to read
% Section 3.1 and the pseudocode; subsystem symbols live with their sections
% (tab:notation-scaffold, tab:notation-ism) and the RNG substream key moved to
% the supplement's reproducibility section. At ~19 rows this fits a normal [t]
% float at \footnotesize, so the Phase-8 [H] override that used to wrap this
% \input at proposed_algorithm.tex is no longer needed.
%
% Kept as a single tabular (not longtable) so the Word build converts it to a
% native w:tbl with a Table caption/SEQ field.
% ============================================================================
\begin{table}[t]
\caption{Core notation: problem, population, and base-GSK symbols used in the
pseudocode (Algorithm~\ref{alg:dt-gsk}) and the equation registry
(Eqs.~\ref{eq:junior-idx}--\ref{eq:rng-substreams}). Subsystem symbols are
defined with their sections (Tables~\ref{tab:notation-scaffold}
and~\ref{tab:notation-ism}); per-equation GSK trial/mask symbols ($u_i$, $v_i$,
$m_{ij}$) are defined in place, and constants take the frozen-configuration
values (Table~\ref{tab:parameters}).}
\label{tab:notation}
\centering
\footnotesize
\begin{tabular}{@{}l@{\hspace{1.2em}}p{0.62\textwidth}@{}}
\toprule
\textbf{Symbol} & \textbf{Meaning}\\
\midrule
\multicolumn{2}{l}{\emph{Problem and population}}\\
\midrule
$f:\mathbb{R}^D\to\mathbb{R}$ & objective (minimized)\\
$D$ & dimension\\
$[\ell,u]^D$ & box bounds ($\ell,u$ scalar, or per-coordinate $\ell_j,u_j$)\\
$\text{MaxFES}$ & evaluation budget (suite-dependent; Section~\ref{sec:alg:complexity})\\
$t$ & evaluations used so far\\
$x = t/\text{MaxFES}\in[0,1]$ & budget fraction\\
$NP_{\text{init}} = 5D$ & initial population size\\
$N_{\min}$ & NLPSR floor (12 bare; 25 at $D\ge50$)\\
$NP(x)$ & current population size (NLPSR schedule)\\
$P=\{x_i\}_{i=1}^{NP}$ & population (rows $x_i\in\mathbb{R}^D$)\\
$x^{*}, f^{*}$ & working incumbent and its value\\
$x^{gb}, f^{gb}$ & global-best (preserved across restarts)\\
\midrule
\multicolumn{2}{l}{\emph{GSK operator}}\\
\midrule
$R_1,R_2,R_3$ & knowledge-source indices (junior: rank-neighbors;\\
              & \quad senior: top/mid/worst groups)\\
$s_J,s_S$ & junior/senior comparison sign (Eq.~\ref{eq:gsk-update}): $+1$ when
the compared source ($x_{R_3}$ junior, $x_{R_2}$ senior) is fitter than $x_i$,
else $-1$\\
$KF$ & knowledge factor (step scale)\\
$KR$ & knowledge ratio (per-dimension update probability)\\
$K_{\exp}$ & junior$\to$senior dimension-schedule exponent\\
$D_{\text{jun}}=\lceil D(1-x)^{K_{\exp}}\rceil$ & \emph{expected} junior-phase coordinate count; a coordinate joins the junior phase with marginal probability $p_{\text{jun}}=D_{\text{jun}}/D$, drawn independently per coordinate on the scalar-mask rows and per block on the linkage-blockwise rows\\
$p$ & senior partition fraction (group size $\operatorname{round}(p\,NP)$)\\
\bottomrule
\end{tabular}
\end{table}
